\scanpage{97}

Écrivons enfin
\[
x'\wedge y'/2q(x')=x\wedge y/2q(x)
\]
On a
\[
\left\{\begin{aligned}x'\wedge y'&=(ad+bc)x\wedge y=\lambda(a^2+b^2)x\wedge y\\q(x')&=(a^2+b^2)q(x)\end{aligned}\right.
\]
donc $x'\wedge y'/q(x')=\lambda x\wedge y/q(x)$ donc la condition est $\lambda=1$. Donc
\[
\unsure{\xi'|\zeta'}=\alpha\xi,\qquad\alpha=a+ib,\quad a^2+b^2\in\Gamma G_m,
\]
donc on a \uline{mono}.

b) C'est un \uline{épi}, en d'autres termes \deleted{\unsure{loc. pour}} $\widetilde f$ est un épi, ou encore \unsure{tt|$\forall$} $\deleted{t}\unsure{z|t}\in\Gamma\Sigma$ se met loc. sous la forme $\widetilde f(\xi)$. Or on a $\varphi(\deleted{t}z,\deleted{t}z)=2q(\deleted{t}z)\in\Gamma G_m$, donc $\deleted{t}z^\perp$ définit un sous-fibré de
$E$ de rang $2$, \deleted{donc} supplémentaire\inserted{$F$} de $\unsure{\mathcal O_S|\mathcal E_S}\cdot\deleted{t}z$, \deleted{\unsure{dont|tout}} et la forme induite par $q$ y est non dég. Comme on est en car. $\ne2$, loc. pour la top. étale on trouve une base orthonormale $x,y$ de $F$.
\[
q(x)=q(y)=1,\quad\varphi(x,y)=0
\]
% archive sheet number: 96
